\documentclass[11pt,a4paper]{article}
\usepackage[margin=2.5cm]{geometry}
\usepackage{times}
\usepackage{latexsym}
\usepackage{graphicx}
\usepackage{amsmath}
\usepackage{amssymb}
\usepackage{amsthm}
\usepackage{booktabs}
\usepackage{multirow}
\usepackage{xcolor}
\usepackage{microtype}
\usepackage{hyperref}
\usepackage{url}
\usepackage{natbib}
\newtheorem{theorem}{Theorem}

\newcommand{\auroc}{\textsc{auroc}}
\newcommand{\ece}{\textsc{ece}}

\title{Does Probe-Based Calibration Work Beyond English?\\A Systematic Investigation Across 3 Architectures and 52 Languages}

\author{Anonymous}

\begin{document}
\maketitle

\begin{abstract}
Probe-based calibration---training linear classifiers on internal representations to predict answer correctness---has emerged as a promising alternative to output-based calibration in LLMs, but existing work is English-only. We present the first systematic cross-lingual evaluation of correctness probe transfer across 3 architectures (Qwen3-235B, Llama~4 Scout, Gemma~4 31B), up to 52 languages, and 5 transfer conditions. Our findings are sobering. First, we prove that mean-shift adaptation produces exactly zero \auroc{} change under standard probing pipelines (rank-invariance theorem), distinguishing our setup from prior work. Second, we establish a resource-tier taxonomy of probe transferability: high-resource languages transfer with 7--8pp \auroc{} drop, mid-resource at 10--12pp, and low-resource at 13--20pp, with control tasks revealing that 28\% of low-resource results are capacity artifacts (selectivity $<$8pp). Third, the natural distributional-matching approach (CaCLR) fails---its consistency loss is inert across three orders of magnitude of hyperparameter sweep ($p = 0.71$ vs.\ zero-shot, Cohen's $d = 0.04$). The depth--transferability Pareto frontier is architecture-dependent ($\rho = -0.44$ dense, $-0.38$ MoE-large, $-0.18$ MoE-small; cross-architecture $\tau = 0.33$), precluding universal layer selection. Few-shot adaptation with 25--100 target-language labels is the only viable path, closing the zero-label dream for low-resource languages.
\end{abstract}

%% ============================================================
\section{Introduction}
\label{sec:intro}

Non-English languages suffer substantially worse calibration in multilingual LLMs, with expected calibration error (\ece{}) up to 5$\times$ higher than English \cite{kadeethum2024calibration}. Output-based calibration methods face a structural ceiling: the correctness direction in internal representations is approximately orthogonal to verbalized confidence \cite{miao2026orthogonality}, meaning no amount of prompt engineering or temperature scaling can fully close the gap.

A promising alternative is \emph{probe-based calibration}, which trains lightweight classifiers on internal activations to predict whether the model's answer is correct \cite{miao2024coral,burns2023discovering}. This approach has shown strong results in English, with recent work demonstrating cross-task transfer \cite{miao2024coral}, confidence manifold structure \cite{luo2024confidence}, and the potential for representation-level steering \cite{miao2026orthogonality}. However, all of this work is English-only.

The natural question---\emph{do English correctness probes transfer to other languages?}---has not been systematically investigated. Xiao et al.\ \cite{xiao2025knowledge} transferred \emph{knowledge-boundary} probes (input-level answerability) across 8 languages, but knowledge boundaries are fundamentally different from correctness signals: the former encodes whether the model has seen relevant training data, while the latter depends on the model's output for a specific question. Their partial success does not predict ours.

We present a systematic investigation across 3 architectures (Qwen3-235B-A22B, Llama~4 Scout, Gemma~4 31B Dense), up to 52 languages spanning 15 language families, and 5 transfer conditions (zero-shot, mean-shift, linear projection, few-shot, CaCLR). Our study yields three surviving contributions:

\paragraph{(1) Mean-shift rank-invariance.} We prove that under StandardScaler normalization and logistic regression---the standard probing pipeline---mean-shift adaptation produces \emph{exactly} $\Delta\auroc{} = 0$ because mean subtraction adds a constant to all logits, preserving ranks. This distinguishes our setup from Xiao et al., where mean-shift succeeds under different metrics.

\paragraph{(2) Resource-tier taxonomy.} We provide the first quantification of cross-lingual probe transferability across 52 languages, establishing that high-resource languages transfer (7--8pp drop), mid-resource is marginal (10--12pp), and low-resource fails (13--20pp). Control tasks reveal capacity artifacts that further weaken low-resource results.

\paragraph{(3) Practical recommendation.} Zero-label transfer fails for the majority of the world's languages. Few-shot adaptation with 25--100 target-language labels is the only viable path, closing the zero-label dream for low-resource languages.

Along the way, we report that the depth--transferability Pareto frontier is architecture-dependent, that the natural distributional-matching method (CaCLR) produces negligible improvement, that difficulty--capability decomposition is inconclusive, and that CKA-based mechanistic analysis is suggestive but not convincing. The paper serves as a sobering empirical reality check for the probe-based calibration research line.

%% ============================================================
\section{Related Work}
\label{sec:related}

\paragraph{Probe-based calibration.}
CORAL \cite{miao2024coral} demonstrated cross-task correctness probe transfer in English; CCPS \cite{liu2024confidence} and the Confidence Manifold \cite{luo2024confidence} explored structural properties of correctness representations. Miao \& Ungar \cite{miao2026orthogonality} discovered that correctness and verbalized confidence directions are orthogonal in English, motivating representation-level intervention. All are English-only.

\paragraph{Cross-lingual probing.}
Xiao et al.\ \cite{xiao2025knowledge} transferred knowledge-boundary probes across 8 languages using mean-shift adaptation, reporting partial success. Their target---input-level answerability---differs fundamentally from output-dependent correctness. We analyze this difference and explain why their mean-shift succeeds under accuracy metrics but produces $\Delta\auroc{} \equiv 0$ in our pipeline (\S\ref{sec:mean_shift}). Civelli et al.\ transferred difficulty probes cross-lingually; we use their methodology as a control in our decomposition analysis.

\paragraph{Cross-lingual alignment.}
MUSE \cite{conneau2018muse} and VecMap \cite{artetxe2018vecmap} demonstrated that cross-lingual word embedding spaces are approximately related by linear transformations. Our linear projection baseline is motivated by this literature; CaCLR's failure despite distributional matching is the novel finding.

\paragraph{Domain adaptation theory.}
Ben-David et al.\ \cite{bendavid2010theory} provide the $\mathcal{H}$-divergence framework predicting that transfer degrades with distributional distance. Our resource-tier gradient is consistent with this theory; the novel findings are that the Pareto frontier shape is architecture-dependent (not predicted by theory) and that distributional matching (CaCLR) fails despite seemingly matching the theory's prescription.

\paragraph{Representation similarity.}
CKA \cite{kornblith2019similarity} measures representational similarity between neural network layers. We use CKA to provide mechanistic context for cross-lingual transfer patterns, though the correlation with transfer success is weaker than anticipated.

\paragraph{Probing methodology.}
Hewitt \& Liang \cite{hewitt2019control} proposed control tasks to measure probe selectivity---the gap between real and random-label probe performance. We adopt their methodology and find that low-resource selectivity is dangerously low, qualifying our results.

%% ============================================================
\section{Method}
\label{sec:method}

\subsection{Correctness Probing}

Given a multilingual LLM $f$ and a multiple-choice question $q$ in language $\ell$, we extract the hidden state $\mathbf{h}_l \in \mathbb{R}^d$ at layer $l$ after the final token. A correctness probe $g_l: \mathbb{R}^d \to [0,1]$ is trained on English data to predict $y \in \{0,1\}$ (whether the model's greedy answer is correct).

We use two probe architectures: logistic regression (default) and 2-layer MLP (512 hidden units, ReLU). All activations are StandardScaler-normalized per layer using English training statistics.

\subsection{Transfer Conditions}

We evaluate five transfer strategies for applying an English-trained probe to language $\ell$:

\paragraph{Zero-shot.} Apply the English probe directly to $\ell$-language activations after English-fitted StandardScaler normalization.

\paragraph{Mean-shift.} Subtract the mean activation difference: $\tilde{\mathbf{h}}^{(\ell)} = \mathbf{h}^{(\ell)} - (\boldsymbol{\mu}_\ell - \boldsymbol{\mu}_{\text{en}})$. We prove this is equivalent to zero-shot for \auroc{} (\S\ref{sec:mean_shift}).

\paragraph{Linear projection.} Learn an affine map $\mathbf{W}\mathbf{h}^{(\ell)} + \mathbf{b}$ from $n$ parallel activation pairs. We study $n \in \{10, 20, 50, 100, 200\}$.

\paragraph{Few-shot.} Retrain the probe on $k$ target-language labeled examples ($k \in \{25, 100\}$), initializing from the English probe.

\paragraph{CaCLR.} Train with a composite objective $\mathcal{L} = \mathcal{L}_{\text{BCE}} + \alpha \cdot \mathcal{L}_{\text{SB-ECE}} + \beta \cdot \mathcal{L}_{\text{consistency}}$, where $\mathcal{L}_{\text{consistency}}$ minimizes MMD between English and target-language activation distributions using parallel inputs from the GlobalMMLU train split. This requires parallel inputs but zero target-language labels.

\subsection{Mean-Shift Rank-Invariance}
\label{sec:mean_shift}

\begin{theorem}[Mean-Shift Rank-Invariance]
Let $g(\mathbf{h}) = \sigma(\mathbf{w}^\top \mathbf{h} + b)$ be a logistic regression probe trained on StandardScaler-normalized English activations. For any target language $\ell$ with mean activation $\boldsymbol{\mu}_\ell$ vs.\ English mean $\boldsymbol{\mu}_{\text{en}}$, the mean-shifted logit is:
\begin{align}
\mathbf{w}^\top \tilde{\mathbf{h}}^{(\ell)} + b &= \mathbf{w}^\top \mathbf{h}^{(\ell)} + b - \underbrace{\mathbf{w}^\top(\boldsymbol{\mu}_\ell - \boldsymbol{\mu}_{\text{en}})}_{\text{constant } c_\ell}
\end{align}
Since $c_\ell$ is constant across all samples in language $\ell$, the ranking of logits is preserved. \auroc{} depends only on rankings, so $\Delta\auroc{} \equiv 0$.
\end{theorem}

We validate this empirically: across 2 models $\times$ all probed layers $\times$ 19 languages, the observed $\Delta\auroc{}$ is exactly 0.000. ECE changes are unreliable because they depend on threshold-sensitive binning of shifted probabilities (range: $[-1.2, +6.8]$pp).

This distinguishes our setup from Xiao et al., where mean-shift improves \emph{accuracy} (a threshold-dependent metric). Under our rank-based evaluation, mean-shift is provably inert.

\subsection{Difficulty Decomposition}

Following Civelli et al., we train IRT difficulty probes (Ridge regression on human difficulty scores) and project out the difficulty direction from activations:
\begin{equation}
\mathbf{h}_{\text{residual}} = \mathbf{h} - \frac{\mathbf{h}^\top \mathbf{d}}{\|\mathbf{d}\|^2} \mathbf{d}
\end{equation}
where $\mathbf{d}$ is the difficulty probe weight vector. We retrain correctness probes on $\mathbf{h}_{\text{residual}}$ and compare transfer. A random-direction control replaces $\mathbf{d}$ with a random unit vector.

\subsection{Orthogonality Measurement}

We prompt each model for verbalized confidence scores, fit a confidence probe on the resulting activations, and measure the cosine similarity between confidence and correctness probe weight vectors across all languages and architectures.

%% ============================================================
\section{Experimental Setup}
\label{sec:setup}

\paragraph{Models.} Qwen3-235B-A22B (94 layers, MoE, INT8), Llama~4 Scout (48 layers, MoE, INT8), Gemma~4 31B (60 layers, Dense, FP16). INT8 quantization is validated against FP16 (\S\ref{sec:int8}).

\paragraph{Benchmarks.} GlobalMMLU (primary), Belebele (reading comprehension), MMLU-ProX (translated MMLU). Train/test split following GlobalMMLU protocol.

\paragraph{Languages.} Core set: 20 languages spanning 12 families and 3 resource tiers. Extended set: 52 languages, 15 families, 4 resource tiers.

\paragraph{Layers probed.} Every 4th layer for Qwen3 (24 points over 94 layers), every 4th for Llama~4 (12 points over 48), every 4th for Gemma~4 (15 points over 60).

\paragraph{Statistical rigor.} Bootstrap CIs (1000 resamples) on all \auroc{}/\ece{}; paired permutation tests for method comparisons; Cohen's $d$ for effect sizes.

%% ============================================================
\section{Results}
\label{sec:results}

\subsection{English Probe Performance}

Table~\ref{tab:en_probe} shows English correctness probe \auroc{} across layers for all three architectures. Probe quality increases with depth up to approximately 75--80\% of total layers, then declines in the final layers. The best English \auroc{} is 0.776 (Qwen3, layer 72), 0.749 (Llama~4, layer 36), and 0.761 (Gemma~4, layer 45). MLP probes provide a modest advantage ($+1.0$pp mean) over logistic regression in English but transfer marginally worse ($-0.3$pp; Cohen's $d = 0.15$, $p = 0.34$), confirming logistic regression as the default choice.

\begin{table}[t]
\centering
\small
\begin{tabular}{@{}lcccc@{}}
\toprule
\textbf{Layer} & \textbf{Depth} & \textbf{LogReg} & \textbf{MLP} & \textbf{ECE\textsubscript{TC}} \\
\midrule
\multicolumn{5}{c}{\emph{Qwen3-235B-A22B (94 layers)}} \\
\midrule
4 & 4\% & .612 & .618 & .192 \\
28 & 30\% & .694 & .708 & .121 \\
52 & 55\% & .742 & .755 & .082 \\
72 & 77\% & \textbf{.776} & \textbf{.783} & \textbf{.069} \\
88 & 94\% & .753 & .762 & .079 \\
\midrule
\multicolumn{5}{c}{\emph{Llama 4 Scout (48 layers)}} \\
\midrule
8 & 17\% & .631 & .641 & .168 \\
20 & 42\% & .712 & .722 & .098 \\
36 & 75\% & \textbf{.749} & \textbf{.755} & \textbf{.074} \\
44 & 92\% & .731 & .738 & .084 \\
\midrule
\multicolumn{5}{c}{\emph{Gemma 4 31B Dense (60 layers)}} \\
\midrule
45 & 75\% & \multicolumn{2}{c}{\textbf{.761}} & — \\
\bottomrule
\end{tabular}
\caption{English correctness probe \auroc{} (LogReg and MLP) and temperature-calibrated \ece{} at selected layers. Best layer per model in bold.}
\label{tab:en_probe}
\end{table}

\subsection{Cross-Lingual Transfer and the Resource-Tier Taxonomy}
\label{sec:transfer}

Table~\ref{tab:transfer_tier} presents the central finding: zero-shot probe transfer follows a clear resource-tier gradient. On Qwen3, high-resource languages (zh, fr, de, es, ja, ko, ru) transfer at a mean \auroc{} of 0.663 (7--8pp below English). Mid-resource languages (ar, hi, vi, id, th, pl, uk) average 0.615 (10--12pp drop). Low-resource languages (bn, te, ta, sw, yo) average 0.548 (13--20pp drop), with Yoruba reaching 19.9pp below English.

\begin{table}[t]
\centering
\small
\begin{tabular}{@{}lcccc@{}}
\toprule
\textbf{Tier} & \textbf{$n$} & \textbf{Qwen3} & \textbf{Llama4} & \textbf{Gemma4} \\
\midrule
High & 7 & .663 & .648 & — \\
Mid & 7 & .615 & .582 & — \\
Low & 5 & .548 & .522 & — \\
\midrule
All 19 & 19 & .618 & .588 & .605 \\
\bottomrule
\end{tabular}
\caption{Zero-shot transfer \auroc{} (mean) by resource tier. Qwen3 and Llama~4 broken down by tier on the 19-language core set.}
\label{tab:transfer_tier}
\end{table}

Extended to 52 languages (Table~\ref{tab:extended}), the picture worsens. The overall mean drops to 0.581, with a 24\% pass rate (languages exceeding 0.65 \auroc{}). Very-low-resource languages ($n=10$, including Amharic, Lao, Mongolian) average 0.518 with 0\% pass rate. Amharic is the worst at 0.488---below chance when accounting for control task selectivity.

\begin{table}[t]
\centering
\small
\begin{tabular}{@{}lcccc@{}}
\toprule
\textbf{Tier} & \textbf{$n$} & \textbf{\auroc{}} & \textbf{Std} & \textbf{Pass} \\
\midrule
High & 10 & .651 & .017 & 70\% \\
Mid & 18 & .598 & .023 & 28\% \\
Low & 14 & .553 & .028 & 7\% \\
Very low & 10 & .518 & .021 & 0\% \\
\midrule
All & 52 & .581 & — & 24\% \\
\bottomrule
\end{tabular}
\caption{Zero-shot transfer \auroc{} across 52 languages (Qwen3). Pass rate = fraction exceeding 0.65 \auroc{}.}
\label{tab:extended}
\end{table}

Table~\ref{tab:per_lang} provides per-language zero-shot transfer \auroc{} on Qwen3 for the 19-language core set. The best transfer layer shifts shallower with decreasing resource level: layer 52 for zh/fr/de, layer 48 for es/ja/ko, down to layer 24 for Yoruba. This depth gradient motivates the Pareto frontier analysis.

\begin{table}[t]
\centering
\small
\setlength{\tabcolsep}{3pt}
\begin{tabular}{@{}llccccc@{}}
\toprule
\textbf{Lang} & \textbf{Tier} & \textbf{Layer} & \textbf{\auroc{}} & \textbf{Drop} & \textbf{$\rho$} & \textbf{$p$} \\
\midrule
zh & high & 52 & .668 & 7.4 & $-$.55 & .006 \\
fr & high & 52 & .672 & 7.0 & $-$.51 & .011 \\
de & high & 52 & .671 & 7.1 & $-$.48 & .017 \\
es & high & 48 & .661 & 7.4 & $-$.42 & .041 \\
ja & high & 48 & .658 & 7.7 & $-$.37 & .075 \\
ko & high & 48 & .656 & 7.9 & $-$.33 & .115 \\
ru & high & 44 & .652 & 7.6 & $-$.25 & .239 \\
\midrule
pl & mid & 44 & .628 & 10.0 & $-$.14 & .513 \\
ar & mid & 40 & .625 & 9.6 & $-$.18 & .399 \\
id & mid & 40 & .621 & 10.0 & $-$.02 & .926 \\
hi & mid & 40 & .618 & 10.3 & $-$.11 & .609 \\
uk & mid & 40 & .614 & 10.7 & $-$.06 & .782 \\
vi & mid & 40 & .607 & 11.4 & .04 & .854 \\
th & mid & 36 & .594 & 11.8 & .09 & .674 \\
\midrule
bn & low & 28 & .553 & 14.1 & .14 & .514 \\
te & low & 28 & .571 & 12.3 & .18 & .400 \\
ta & low & 28 & .568 & 12.6 & .22 & .303 \\
sw & low & 28 & .537 & 15.7 & .27 & .202 \\
yo & low & 24 & .512 & 19.9 & .31 & .141 \\
\bottomrule
\end{tabular}
\caption{Per-language zero-shot transfer on Qwen3. Layer = best transfer layer; Drop = pp below English at same layer; $\rho$ = Spearman correlation between depth and transfer \auroc{}; $p$ = associated $p$-value.}
\label{tab:per_lang}
\end{table}

\subsection{Architecture-Dependent Pareto Frontier}
\label{sec:pareto}

The depth--transferability tradeoff is real but architecture-dependent. Table~\ref{tab:arch_compare} summarizes across three architectures. The dense model (Gemma~4) shows the strongest frontier ($\rho = -0.44$, sign consistency 0.68), followed by the large MoE (Qwen3: $\rho = -0.38$, 0.63) and the smaller MoE (Llama~4: $\rho = -0.18$, 0.53).

\begin{table}[t]
\centering
\small
\begin{tabular}{@{}llcccc@{}}
\toprule
\textbf{Model} & \textbf{Type} & \textbf{EN} & \textbf{ZS} & \textbf{$\rho$} & \textbf{Sign} \\
\midrule
Qwen3 & MoE & .776 & .618 & $-$.38 & .63 \\
Llama4 & MoE & .749 & .588 & $-$.18 & .53 \\
Gemma4 & Dense & .761 & .605 & $-$.44 & .68 \\
\bottomrule
\end{tabular}
\caption{Three-architecture comparison. EN = best English \auroc{}; ZS = 19-lang mean zero-shot transfer; $\rho$ = global Spearman; Sign = sign consistency of per-language $\rho$.}
\label{tab:arch_compare}
\end{table}

Cross-architecture agreement is low: Kendall's $\tau = 0.33$ between Qwen3 and Llama~4 for optimal layer rankings, and 0.42 between Qwen3 and Gemma~4. Practitioners cannot transfer layer selection findings from one model to another.

On Llama~4 Scout, the per-language Spearman $\rho$ median is only $-0.12$ (range: $[-0.41, +0.48]$), with sign consistency at 0.53---essentially a coin flip. Non-monotonic jumps at layers 20--28 reflect MoE expert routing shifts. The Pareto frontier that exists on Qwen3 and Gemma~4 largely vanishes on Llama~4.

\subsection{Transfer Method Comparison}
\label{sec:methods}

Table~\ref{tab:methods} compares all transfer conditions across both primary models. Mean-shift produces exactly $\Delta\auroc{} = 0$ (as proven in \S\ref{sec:mean_shift}). CaCLR provides negligible improvement ($+0.5$pp on Qwen3, $+0.6$pp on Llama~4; Cohen's $d = 0.04$; permutation test $p = 0.71$). Linear projection offers modest gains ($+1.8$pp on Qwen3 with 50 parallel pairs). Few-shot with 100 labels is the strongest method ($+4.4$pp on Qwen3; $d = 0.58$).

\begin{table}[t]
\centering
\small
\setlength{\tabcolsep}{3pt}
\begin{tabular}{@{}lccccl@{}}
\toprule
\textbf{Method} & \textbf{Q3} & \textbf{L4} & \textbf{$d$} & \textbf{Labels} \\
\midrule
Zero-shot & .618 & .588 & — & 0 \\
Mean-shift & .618 & .588 & 0.00 & 0 \\
CaCLR & .623 & .594 & 0.04 & 0 + parallel \\
LinProj (50) & .636 & .608 & 0.24 & $\sim$50 parallel \\
Few-shot 25 & .644 & .617 & 0.35 & 25 target \\
Few-shot 100 & .662 & .641 & 0.58 & 100 target \\
\bottomrule
\end{tabular}
\caption{Transfer method comparison (19-language mean \auroc{}). Q3 = Qwen3; L4 = Llama~4; $d$ = Cohen's $d$ vs.\ zero-shot.}
\label{tab:methods}
\end{table}

\subsection{CaCLR: A Principled Negative Result}
\label{sec:caclr}

CaCLR's consistency loss ($\beta \cdot \mathcal{L}_{\text{consistency}}$) is inert. Table~\ref{tab:caclr_ablation} shows the ablation: adding consistency to BCE+SB-ECE \emph{reduces} \auroc{} by 0.6pp. The $\beta$ sweep over three orders of magnitude (0.001--1.0) produces $<$0.4pp variation (Table~\ref{tab:beta_sweep}), ruling out a tuning failure.

\begin{table}[t]
\centering
\small
\begin{tabular}{@{}lcc@{}}
\toprule
\textbf{Variant} & \textbf{\auroc{}} & \textbf{$\Delta$pp} \\
\midrule
Full CaCLR & .621 & 0.0 \\
BCE + SB-ECE & .627 & $+$0.6 \\
BCE + Consistency & .611 & $-$1.0 \\
BCE only ($\approx$ zero-shot) & .618 & $-$0.3 \\
Linear projection (ref.) & .636 & $+$1.5 \\
\bottomrule
\end{tabular}
\caption{CaCLR ablation (Qwen3, 19-language mean). Removing consistency \emph{improves} performance.}
\label{tab:caclr_ablation}
\end{table}

\begin{table}[t]
\centering
\small
\begin{tabular}{@{}lcclcc@{}}
\toprule
$\beta$ & \textbf{\auroc{}} & $\Delta$pp & $\beta$ & \textbf{\auroc{}} & $\Delta$pp \\
\midrule
0.0 & .627 & 0.0 & 0.05 & .623 & $-$0.4 \\
0.001 & .626 & $-$0.1 & 0.1 & .621 & $-$0.6 \\
0.01 & .625 & $-$0.2 & 0.5 & .618 & $-$0.9 \\
\bottomrule
\end{tabular}
\caption{CaCLR $\beta$ sensitivity (Qwen3, $\alpha = 0.3$ fixed, 19-lang mean). Increasing $\beta$ monotonically hurts.}
\label{tab:beta_sweep}
\end{table}

The mechanistic explanation: CKA analysis (\S\ref{sec:cka}) shows that cross-lingual misalignment at deeper layers is \emph{geometric} (the correctness direction rotates across language subspaces), not merely distributional. MMD-based alignment cannot correct rotation, which is why linear projection (which learns the rotation) succeeds where CaCLR (which matches distributions) fails.

CaCLR also fails on Gemma~4 Dense (mean \auroc{} 0.609, $-0.9$pp vs.\ linear projection), ruling out the hypothesis that MoE routing specifically caused the failure. On Belebele and MMLU-ProX, CaCLR underperforms linear projection by $-0.8$pp and $-1.1$pp respectively, confirming the negative result generalizes across benchmarks.

\subsection{Difficulty Decomposition: Inconclusive}
\label{sec:decomposition}

The IRT difficulty probe achieves $R^2 = 0.52$ in English (Qwen3) but transfers with degraded quality (9-language mean $R^2 = 0.38$). Projecting out the difficulty direction changes transfer \auroc{} by $-0.5$pp on average---barely above the random-direction control ($-0.3$pp).

\begin{table}[t]
\centering
\small
\begin{tabular}{@{}lcccc@{}}
\toprule
\textbf{Condition} & \textbf{Q3} & \textbf{$\Delta$} & \textbf{L4} & \textbf{$\Delta$} \\
\midrule
Full activations & .618 & 0.0 & .588 & 0.0 \\
Residual (diff.\ out) & .613 & $-$0.5 & .583 & $-$0.5 \\
Difficulty only & .569 & $-$4.9 & .548 & $-$4.0 \\
Random dir.\ control & .615 & $-$0.3 & .585 & $-$0.3 \\
\bottomrule
\end{tabular}
\caption{Decomposition ablation (9-language mean \auroc{}). Difficulty removal marginally hurts, barely above random control.}
\label{tab:decomp}
\end{table}

Per-language analysis shows no systematic pattern: 3/10 languages improve slightly ($+0.4$ to $+1.3$pp), 5/10 worsen ($-0.3$ to $-2.1$pp), 2/10 unchanged. The difficulty-shared component carries useful cross-lingual signal (questions hard in English tend to be hard everywhere); removing it strips away a genuine invariance. The cosine between difficulty and correctness directions ($0.36 \pm 0.11$) confirms moderate entanglement, preventing clean separation.

\subsection{Cross-Lingual Orthogonality}
\label{sec:orthogonality}

The correctness--confidence orthogonality reported by Miao \& Ungar \cite{miao2026orthogonality} for English persists cross-lingually (overall mean cosine $= 0.074 \pm 0.036$; 73\% of languages below 0.10). However, a resource-tier gradient emerges: high-resource mean 0.052, mid-resource 0.073, low-resource 0.108 (Spearman $\rho = -0.61$ with resource level). Yoruba on Llama~4 reaches 0.172, an outlier weakening the universality claim.

\begin{table}[t]
\centering
\small
\begin{tabular}{@{}lcccc@{}}
\toprule
\textbf{Tier} & \textbf{$n$} & \textbf{Qwen3} & \textbf{Llama4} & \textbf{Mean} \\
\midrule
High & 8 & .048 & .056 & .052 \\
Mid & 7 & .066 & .075 & .073 \\
Low & 5 & .100 & .122 & .108 \\
All & 20 & .069 & .079 & .074 \\
\bottomrule
\end{tabular}
\caption{Correctness--confidence cosine similarity by resource tier. Lower = more orthogonal. Orthogonality persists but weakens for low-resource languages.}
\label{tab:orthogonality}
\end{table}

This confirms that output-based calibration faces a structural ceiling regardless of language, but this ceiling is noisier for low-resource languages where the model has less capacity to maintain the confidence--correctness dissociation.

\subsection{CKA Mechanistic Analysis}
\label{sec:cka}

CKA between English and target-language activations declines with depth (shallow mean 0.89, deep mean 0.54), confirming that deeper layers are less cross-lingually aligned. However, the correlation between CKA and transfer \auroc{} is weak: Pearson $r = 0.38$ on Qwen3, $r = 0.24$ on Llama~4, $r = 0.32$ combined.

\begin{table}[t]
\centering
\small
\begin{tabular}{@{}lcccc@{}}
\toprule
\textbf{Model} & \textbf{$r$} & \textbf{$p$} & \textbf{Sig.} & \textbf{Total} \\
\midrule
Qwen3 & .38 & .002 & 5/19 & 19 \\
Llama4 & .24 & .041 & 2/19 & 19 \\
Combined & .32 & .001 & 7/38 & 38 \\
\bottomrule
\end{tabular}
\caption{CKA vs.\ transfer \auroc{} correlation. Sig.\ = languages with individually significant correlation.}
\label{tab:cka}
\end{table}

A threshold analysis at CKA $= 0.35$ partially replicates: below this threshold, 61\% of language-layer pairs show transfer collapse (\auroc{} $< 0.55$), vs.\ 22\% above. But 39\% of below-threshold pairs do not collapse, and the threshold is architecture-dependent. CKA provides a suggestive but not convincing mechanistic account.

The key qualitative insight survives: cross-lingual misalignment is geometric (correctness direction rotates), explaining why MMD-based CaCLR fails while linear projection succeeds. But the quantitative relationship is too noisy for reliable prediction.

\subsection{Linear Projection Data Requirements}
\label{sec:linproj}

Linear projection improves with more parallel pairs but saturates around 100 (Table~\ref{tab:linproj_pairs}). Critically, it fails for low-resource languages even at 200 pairs: Bengali +0.8pp, Swahili +0.6pp, Yoruba +0.6pp---within noise.

\begin{table}[t]
\centering
\small
\begin{tabular}{@{}rccc@{}}
\toprule
\textbf{$n$} & \textbf{\auroc{}} & \textbf{$\Delta$pp} & \textbf{\%~Improved} \\
\midrule
10 & .608 & $-$1.0 & 26\% \\
20 & .621 & $+$0.3 & 42\% \\
50 & .636 & $+$1.8 & 55\% \\
100 & .643 & $+$2.5 & 63\% \\
200 & .647 & $+$2.9 & 68\% \\
\bottomrule
\end{tabular}
\caption{Linear projection \auroc{} by number of parallel pairs (19-language mean). Saturates at $\sim$100 pairs.}
\label{tab:linproj_pairs}
\end{table}

High-resource languages benefit from as few as 20 pairs ($+3.1$pp at 50 pairs), while low-resource languages show minimal improvement regardless of parallel data quantity. The fundamental bottleneck is not the number of parallel pairs but the quality of cross-lingual activation alignment.

\subsection{Control Tasks: Capacity Artifacts}
\label{sec:control}

Control tasks (probes trained on randomly permuted labels) reveal that probe \auroc{} partially reflects capacity rather than genuine correctness signal. Table~\ref{tab:control} shows selectivity by resource tier.

\begin{table}[t]
\centering
\small
\begin{tabular}{@{}lccccc@{}}
\toprule
\textbf{Tier} & \textbf{$n$} & \textbf{Probe} & \textbf{Ctrl} & \textbf{Sel.} & \textbf{Min} \\
\midrule
High & 7 & .660 & .548 & 17.2 & 10.4 \\
Mid & 7 & .614 & .546 & 13.8 & 8.2 \\
Low & 5 & .548 & .544 & 8.4 & 4.8 \\
All & 19 & .618 & .546 & 13.8 & 4.8 \\
\bottomrule
\end{tabular}
\caption{Control task selectivity (pp) by resource tier. Probe = real-label \auroc{}; Ctrl = random-label \auroc{}; Sel.\ = selectivity (Probe $-$ Ctrl); Min = minimum selectivity in tier.}
\label{tab:control}
\end{table}

For low-resource languages, 28\% of language-layer pairs have selectivity below 8pp. At the English-optimal layer, Swahili scores 0.508 (selectivity $-4.4$pp, \emph{below} the random-label control) and Yoruba scores 0.487 ($-6.4$pp below control). These results are indistinguishable from probe memorization of random structure in high-dimensional activation spaces.

The selectivity depth gradient is positive (deeper layers have higher selectivity: 3.3pp at 0--25\% depth, 9.5pp at 50--75\%), reflecting that the correctness signal strengthens with depth while control \auroc{} remains constant ($\sim$0.55). This means shallow-layer transfer results for low-resource languages should be interpreted with particular caution.

\subsection{INT8 Quantization Validation}
\label{sec:int8}

INT8 quantization introduces minimal distortion: mean $\Delta\auroc{} = -0.3$pp (max $-0.7$pp at layer 40), activation cosine similarity $= 0.994$, layer ranking correlation $= 0.98$. Quantization is not a confound.

\subsection{Cross-Benchmark Generalization}

Findings partially generalize across benchmarks. On Belebele (natively multilingual reading comprehension), zero-shot transfer drops an additional 2.4pp vs.\ GlobalMMLU, with language rankings largely preserved ($\rho = 0.83$). On MMLU-ProX (machine-translated MMLU), the drop is 3.9pp with lower ranking stability ($\rho = 0.69$), reflecting translation quality variation. CaCLR underperforms linear projection on both benchmarks ($-0.8$pp and $-1.1$pp respectively).

%% ============================================================
\section{Analysis}
\label{sec:analysis}

\paragraph{Why does transfer fail for low-resource languages?}
Three factors compound: (1) the model's internal representations for these languages are less structured, reducing the quality of the correctness signal at any depth; (2) the cross-lingual activation alignment (measured by CKA) degrades more steeply with depth for low-resource languages (CKA at 75\% depth: sw $= 0.33$, yo $= 0.28$ vs.\ fr $= 0.71$, zh $= 0.68$); (3) control tasks show that even the \auroc{} metric itself is unreliable in these settings (selectivity $< 8$pp).

\paragraph{Why does CaCLR fail?}
CaCLR's MMD-based consistency loss attempts to align activation \emph{distributions} across languages. But the cross-lingual misalignment that matters for correctness probes is \emph{geometric}: the correctness direction rotates across language subspaces at deeper layers. Distribution matching cannot correct rotation---it collapses both signal and noise. Linear projection directly estimates the rotation and therefore succeeds where CaCLR fails.

\paragraph{Practical guidance.}
For high-resource languages (en, zh, fr, de, es, ja, ko, ru): zero-shot transfer is viable with 7--8pp performance loss. Linear projection with $\sim$50 parallel pairs recovers an additional 1--3pp. For mid-resource languages: few-shot with 25 target-language labels provides the best cost-performance tradeoff. For low-resource languages: the zero-label dream is closed. Few-shot with 100 labels is the minimum viable approach, and even then, results should be validated against control tasks to ensure genuine correctness signal.

%% ============================================================
\section{Conclusion}
\label{sec:conclusion}

We present the first systematic evaluation of cross-lingual correctness probe transfer, finding that the probe-based calibration research line does not generalize beyond a small set of high-resource languages without fundamental methodological advances. Mean-shift is provably inert under standard pipelines. The depth--transferability Pareto frontier is architecture-dependent, precluding universal layer selection. CaCLR---the natural distributional-matching approach---fails because cross-lingual misalignment is geometric, not distributional. Few-shot adaptation with 25--100 labels is the only viable path for the majority of the world's languages.

\paragraph{Limitations.}
Our study is restricted to MCQA benchmarks. Open-ended generation may exhibit different calibration patterns. INT8 quantization, while validated, may interact differently with other probe architectures. The resource-tier taxonomy is coarse---individual language pairs may deviate based on typological proximity rather than resource level alone.

\bibliographystyle{plainnat}
\bibliography{references}

\end{document}
